\documentclass[uplatex,dvipdfmx,a4paper,11pt]{jsarticle}
\input{../../../templates/note.tex}
\noteprint{量子力学1　第2章「波束とシュレーディンガー方程式・定常状態」}{2026-11-25}
%% note-moto	problems\量子力学1\shou02\量子力学1-2.tex
\begin{document}
\begin{kirokuichiran}
  \ichirangyou{\hyperlink{bt-m-009-phase}{1}}{L1・計算}{○△}{C}{2026-10-14 (1)記述 (2)計算}
  \ichirangyou{2}{L1・計算}{}{未着手}{}
  \ichirangyou{\hyperlink{bt-m-011-operator}{3}}{L1・計算}{◎◎◎}{S}{2026-10-14}
  \ichirangyou{4}{L1・計算}{◎◎○}{A}{2026-11-21 (3)記述}
  \ichirangyou{\hyperlink{bt-m-012-order}{5}}{L2・計算}{◎◎◎}{S}{2026-10-19 2回目}
  \ichirangyou{6}{L2・計算}{}{未着手}{}
  \ichirangyou{\hyperlink{bt-m-016-stationary}{7}}{L3・導出}{◎△◎}{B}{2026-10-29 (2)方針}
  \ichirangyou{8}{L3・導出}{}{未着手}{}
  \ichirangyou{9}{L4・導出}{}{未着手}{}
\end{kirokuichiran}

\mondaihen
\level{1}
% toubun: kiroku\toubun\量子力学1\009_2026-10-14.tex
\mondaibangou{1}
\begin{mondai}{009-phase}{計算}{平面波の進む向きと位相速度，進行波と定在波の違い}
\begin{teisuu}
  プランク定数 & $h$ & \qty{6.626e-34}{\joule\second}\\
  換算プランク定数 & $\hbar$ & \qty{1.055e-34}{\joule\second}\\
  電子の質量 & $\melec$ & \qty{9.109e-31}{\kilogram}\\
  電気素量 & $e$ & \qty{1.602e-19}{\coulomb}\\
  真空中の光速 & $c$ & \qty{2.998e8}{\metre\per\second}\\
\end{teisuu}
$k>0$，$\omega>0$ を定数とし，$x$ 軸上の波を表す複素数値の関数を考える。
  \begin{shoumon}
    \item 次の3つの波が，それぞれ $x$ の正・負のどちら向きに進むかを，位相が一定の点の動きから答え，その速さ（位相速度）を求めよ。
    \[ \psi_{a}=\exp[-\ii(kx-\omega t)],\qquad \psi_{b}=\exp[\ii(kx+\omega t)],\qquad \psi_{c}=\exp[-\ii(2kx+\omega t)] \]
    \item 2つの平面波の和 $\psi_{d}=\exp[\ii(kx-\omega t)]+\exp[\ii(-kx-\omega t)]$ を，$x$ だけの関数と $t$ だけの関数の積に書き直せ。また，$\abs{\psi_{d}}=0$ となる点の位置を求め，$\psi_{d}$ が進む波でないことを説明せよ。
  \end{shoumon}
\notekiroku{2026-10-14 C ○△}
\begin{kotaeutsushi}{2026-10-14}
\toi{1}{○}{記述}
\begin{align*}
&\psi_{a}:\ \text{右},\ \frac{\omega t}{k}\qquad\psi_{b}:\ \text{左}\ -\frac{\omega t}{k}\\
&\psi_{c}:\ \text{左}\ -\frac{\omega t}{2k}
\end{align*}
\naze{向きと速さの答えだけで，位相が一定の点の動きの根拠が無い。たとえば $\psi_{c}$ では位相 $-(2kx+\omega t)=$ 一定 から $2k\,\dv{x}{t}+\omega=0$，$\dv{x}{t}=-\omega/(2k)$ である。書いた $\omega t/k$ などは位相が一定の点の位置（$+$定数）で，速さはその係数 $\omega/k$，$\omega/k$，$\omega/(2k)$（正の量）である。}
\toi{2}{△}{計算}
\begin{align*}
&\ee^{\ii kx}\ee^{-\ii\omega t}+\ee^{-\ii kx}\ee^{-\ii\omega t}\\
&\quad=(\ee^{\ii kx}+\ee^{-\ii kx})\ee^{-\ii\omega t}\\
&\quad=\ayamari{2\cos x}{2\cos kx}\,\ee^{-\ii\omega t}
\end{align*}
\naze{$\ee^{\ii kx}+\ee^{-\ii kx}=2\cos kx$ で，$k$ が落ちている。以降はこの誤りを引き継ぐ。}
\begin{align*}
&\abs{\psi_{d}}=\abs{2\cos x}=0\\
&\therefore\ x=\frac{\pi}{2}+n\pi\quad(n=\pm\fumei,2,3,\dots)
\end{align*}
\naze{正しくは $\abs{\psi_{d}}=\abs{2\cos kx}=0$ から $kx=\pi/2+j\pi$，すなわち $x=(2j+1)\pi/(2k)$（$j$ は $0$ を含むすべての整数）である。}
$0$ になる点が $t$ に依存しないため，定常波。
\end{kotaeutsushi}
\begin{kizuki}
\item 位相が $\theta=\alpha x+\beta t+\text{定数}$ の形なら，$\theta$ が一定の点では $\alpha\,\dv{x}{t}+\beta=0$ で $\dv{x}{t}=-\beta/\alpha$ となる。この1行を書いてから，向き（符号）と速さ（正の量）を分けて答える。
\item $\cos$ の引数は無次元なので，$\cos x$ や $x=\pi/2+n\pi$ のように $x$ が $k$ を伴わずに出てきたら $k$ の書き落としを疑う。零点の答えに $k$ が入っているかで検算する。
\end{kizuki}
\begin{tadashii}
\houshin 位相 $\theta(x,t)$ が一定の点が時間とともにどう動くかを追う。位相が $\theta=\alpha x+\beta t+\text{定数}$ の形なら，$\theta$ が一定の点では $\alpha\,\dd x/\dd t+\beta=0$ なので，速度は $-\beta/\alpha$ であり，その符号で向きが決まる。
\keisan
\shoumonkai{1} 各波の位相の $\alpha$，$\beta$ を読み取る。
\begin{align*}
  \psi_{a}&:\ \theta_{a}=-kx+\omega t, & \dv{x}{t}&=-\frac{\omega}{-k}=\frac{\omega}{k}>0,\\
  \psi_{b}&:\ \theta_{b}=kx+\omega t, & \dv{x}{t}&=-\frac{\omega}{k}<0,\\
  \psi_{c}&:\ \theta_{c}=-2kx-\omega t, & \dv{x}{t}&=-\frac{-\omega}{-2k}=-\frac{\omega}{2k}<0.
\end{align*}
よって $\psi_{a}$ は $x$ の正の向きに速さ $\omega/k$，$\psi_{b}$ は負の向きに速さ $\omega/k$，$\psi_{c}$ は負の向きに速さ $\omega/(2k)$ で進む。

\shoumonkai{2} 共通な因子 $\ee^{-\ii\omega t}$ をくくり出し，Euler の公式を使う。
\[ \psi_{d}=\ee^{-\ii\omega t}\left(\ee^{\ii kx}+\ee^{-\ii kx}\right)=2\cos(kx)\,\ee^{-\ii\omega t} \]
$\abs{\ee^{-\ii\omega t}}=1$ なので $\abs{\psi_{d}}=2\abs{\cos kx}$ である。これが0になるのは $kx=(j+1/2)\pi$（$j$ は整数），すなわち
\[ x=\frac{(2j+1)\pi}{2k} \]
のときで，この位置は $t$ によらない。$\abs{\psi_{d}}$ の形全体が時間とともに動かないので，$\psi_{d}$ は進む波ではなく定在波である（右向きと左向きの同じ振幅の波の重ね合わせ）。
\kotae{\shoumonkai{1} $\psi_{a}$：$x$ の正の向き，$\omega/k$；$\psi_{b}$：$x$ の負の向き，$\omega/k$；$\psi_{c}$：$x$ の負の向き，$\omega/(2k)$ \shoumonkai{2} $\psi_{d}=2\cos(kx)\,\ee^{-\ii\omega t}$，$\abs{\psi_{d}}=0$ は $x=(2j+1)\pi/(2k)$（$j$ は整数）で時間によらない（定在波）}
\begin{ronten}
  \item \shoumonkai{1} 位相が一定の点の速度を位相の中の $x$ と $t$ の係数の比から求め，その符号で向きを決めたことを書く（位相全体の符号は向きを変えない）。
  \item \shoumonkai{2} $\abs{\ee^{-\ii\omega t}}=1$ から $\abs{\psi_{d}}$ が $x$ だけで決まり，零点の位置が $t$ によらないことを，進む波でない理由として書く。
\end{ronten}
\kenzan 次元：$\omega/k$ は $\unit{\per\second}/\unit{\per\metre}=\unit{\metre\per\second}$ で速さである。$\psi_{a}$ は右向きの波 $\ee^{\ii(kx-\omega t)}$ の複素共役で，実部はどちらも $\cos(kx-\omega t)$ だから同じ向きに進むことと合う。$\psi_{c}$ の実部 $\cos(2kx+\omega t)=\cos[2k(x+\omega t/(2k))]$ は左向きに速さ $\omega/(2k)$ で平行移動する形であり，結果と一致する。
\tsumazuki 位相全体の符号を反転しても向きは変わらない。向きを決めるのは位相の中の $x$ と $t$ の係数の相対的な符号である。
\end{tadashii}
\end{mondai}

% toubun: kiroku\toubun\量子力学1\011_2026-10-14.tex
\mondaibangou{3}
\begin{mondai}{011-operator}{計算}{運動量演算子 $\hat{p}_{x}$ の作用と固有関数・固有値（$\hat{p}_{x}$ と $\hat{p}_{x}^{2}$ の違い）}
\begin{teisuu}
  プランク定数 & $h$ & \qty{6.626e-34}{\joule\second}\\
  換算プランク定数 & $\hbar$ & \qty{1.055e-34}{\joule\second}\\
  電子の質量 & $\melec$ & \qty{9.109e-31}{\kilogram}\\
  電気素量 & $e$ & \qty{1.602e-19}{\coulomb}\\
  真空中の光速 & $c$ & \qty{2.998e8}{\metre\per\second}\\
\end{teisuu}
1次元の運動量演算子を $\hat{p}_{x}f=-\ii\hbar\dv{f}{x}$ とし，$k>0$ を定数とする。
  \begin{shoumon}
    \item $f_{1}(x)=\ee^{\ii kx}$ に $\hat{p}_{x}$ を作用させ，$f_{1}$ が $\hat{p}_{x}$ の固有関数であることを示して固有値を求めよ。
    \item $f_{2}(x)=\cos kx$ に $\hat{p}_{x}$ を作用させ，$f_{2}$ が $\hat{p}_{x}$ の固有関数でないことを示せ。また，$f_{2}$ を $\hat{p}_{x}$ の固有関数の和で表し，それぞれの固有値を答えよ。
    \item $f_{2}$ に $\hat{p}_{x}^{2}$ を作用させ，$f_{2}$ が $\hat{p}_{x}^{2}$ の固有関数であることを示して固有値を求めよ。さらに，自由粒子のハミルトニアン $\hat{H}=\hat{p}_{x}^{2}/(2m)$ について，$f_{2}$ のエネルギー固有値を求めよ。
  \end{shoumon}
\notekiroku{2026-10-14 S ◎◎◎}
\begin{kotaeutsushi}{2026-10-14}
\toi{1}{◎}{}
\[ \hat{p}_{x}f_{1}=-\ii\hbar\frac{\dd}{\dd x}\ee^{\ii kx}=\hbar k\,\ee^{\ii kx} \]
固有値 $\hbar k$，固有関数 $\ee^{\ii kx}$
\toi{2}{◎}{}
\[ \hat{p}_{x}f_{2}=-\ii\hbar\frac{\dd}{\dd x}\cos kx=\ii\hbar k\sin kx \]
となり，固有関数でない。
\[ f_{2}=\frac{1}{2}(\ee^{\ii kx}+\ee^{-\ii kx}) \]
$\ee^{\ii kx}\cdots\hbar k$，\quad $\ee^{-\ii kx}\cdots-\hbar k$
\toi{3}{◎}{}
\[ \hat{p}_{x}^{2}f_{2}=-\hbar^{2}\frac{\dd^{2}}{\dd x^{2}}\cos kx=\hbar^{2}k^{2}\cos kx \]
固有値 $\hbar^{2}k^{2}$
\[ \hat{H}f_{2}=\frac{\hbar^{2}k^{2}}{2m}\cos kx \]
エネルギー固有値 $(\hbar k)^{2}/(2m)$
\end{kotaeutsushi}
\end{mondai}

\level{2}
% toubun: kiroku\toubun\量子力学1\012_2026-10-19.tex
\mondaibangou{5}
\begin{mondai}{012-order}{計算}{座標と運動量の演算子の順序（交換関係）と，3次元での $\hat{\vect{p}}$ の作用}
\begin{teisuu}
  プランク定数 & $h$ & \qty{6.626e-34}{\joule\second}\\
  換算プランク定数 & $\hbar$ & \qty{1.055e-34}{\joule\second}\\
  電子の質量 & $\melec$ & \qty{9.109e-31}{\kilogram}\\
  電気素量 & $e$ & \qty{1.602e-19}{\coulomb}\\
  真空中の光速 & $c$ & \qty{2.998e8}{\metre\per\second}\\
\end{teisuu}
3次元で，$\hat{\vect{r}}$ は関数に $\vect{r}=(x,y,z)$ を掛け，$\hat{\vect{p}}=-\ii\hbar\nabla$ とする。$f(\vect{r})$ は任意のなめらかな関数とする。
  \begin{shoumon}
    \item $(\hat{x}\hat{p}_{x}-\hat{p}_{x}\hat{x})f$ と $(\hat{x}\hat{p}_{y}-\hat{p}_{y}\hat{x})f$ を計算せよ。
    \item $\hat{\vect{r}}\cdot\hat{\vect{p}}=\hat{x}\hat{p}_{x}+\hat{y}\hat{p}_{y}+\hat{z}\hat{p}_{z}$，$\hat{\vect{p}}\cdot\hat{\vect{r}}=\hat{p}_{x}\hat{x}+\hat{p}_{y}\hat{y}+\hat{p}_{z}\hat{z}$ とする（右側の演算子から先に作用させる）。(1) の結果を用いて，$(\hat{\vect{p}}\cdot\hat{\vect{r}})f-(\hat{\vect{r}}\cdot\hat{\vect{p}})f$ を $f$ で表せ。
    \item $b>0$ を定数，$r=\abs{\vect{r}}$ として $\phi(\vect{r})=\exp[-r^{2}/(2b^{2})]$ とする。$\hat{\vect{p}}\phi$，$(\hat{\vect{r}}\cdot\hat{\vect{p}})\phi$，$(\hat{\vect{p}}\cdot\hat{\vect{r}})\phi$ を計算し，(2) の結果を満たすことを確かめよ。
  \end{shoumon}
\notekiroku{2026-10-14 B ◎△○ → 2026-10-19 S ◎◎◎}
\begin{kotaeutsushi}{2026-10-19}
\toi{1}{◎}{}
\begin{align*}
\hat{x}\hat{p}_{x}f&=x\left(-\ii\hbar\frac{\partial}{\partial x}\right)f=-\ii\hbar x\pdv{f}{x}\\
\hat{p}_{x}\hat{x}f&=-\ii\hbar\frac{\partial}{\partial x}(xf)\\
&=-\ii\hbar\left(1\cdot f+x\pdv{f}{x}\right)\\
&=-\ii\hbar\left(f+x\pdv{f}{x}\right)
\end{align*}
\begin{align*}
&(\hat{x}\hat{p}_{x}-\hat{p}_{x}\hat{x})f\\
&\quad=-\ii\hbar x\pdv{f}{x}-\left(-\ii\hbar\left(f+x\pdv{f}{x}\right)\right)\\
&\quad=\ii\hbar f
\end{align*}
\begin{align*}
\hat{x}\hat{p}_{y}f&=x\left(-\ii\hbar\frac{\partial}{\partial y}\right)f=-\ii\hbar x\pdv{f}{y}\\
\hat{p}_{y}\hat{x}f&=-\ii\hbar\frac{\partial}{\partial y}(xf)=-\ii\hbar\left(0+x\pdv{f}{y}\right)
\end{align*}
\begin{align*}
&(\hat{x}\hat{p}_{y}-\hat{p}_{y}\hat{x})f\\
&\quad=-\ii\hbar x\pdv{f}{y}-\left(-\ii\hbar x\pdv{f}{y}\right)=0
\end{align*}
\toi{2}{◎}{}
\begin{align*}
&(\hat{\vect{p}}\cdot\hat{\vect{r}})f-(\hat{\vect{r}}\cdot\hat{\vect{p}})f\\
&\quad=(\hat{p}_{x}\hat{x}+\hat{p}_{y}\hat{y}+\hat{p}_{z}\hat{z})f\\
&\qquad-(\hat{x}\hat{p}_{x}+\hat{y}\hat{p}_{y}+\hat{z}\hat{p}_{z})f\\
&\quad=-(\hat{x}\hat{p}_{x}-\hat{p}_{x}\hat{x})f-(\hat{y}\hat{p}_{y}-\hat{p}_{y}\hat{y})f\\
&\qquad-(\hat{z}\hat{p}_{z}-\hat{p}_{z}\hat{z})f\\
&\quad=-3\ii\hbar f
\end{align*}
\toi{3}{◎}{}
\[ \hat{\vect{p}}\phi=-\ii\hbar\nabla\phi \]
\begin{align*}
\pdv{\phi}{x}&=\frac{\partial}{\partial x}\left[\exp\left(-\frac{x^{2}+y^{2}+z^{2}}{2b^{2}}\right)\right]\\
&=-\frac{x}{b^{2}}\exp\left(-\frac{r^{2}}{2b^{2}}\right)
\end{align*}
\begin{align*}
-\ii\hbar\nabla\phi&=\frac{\ii\hbar}{b^{2}}\exp\left(-\frac{r^{2}}{2b^{2}}\right)(x,y,z)\\
&=\left(\frac{\ii\hbar}{b^{2}}\exp\left(-\frac{r^{2}}{2b^{2}}\right)\right)\vect{r}
\end{align*}
\begin{align*}
(\hat{\vect{r}}\cdot\hat{\vect{p}})\phi&=\hat{x}\hat{p}_{x}\phi+\hat{y}\hat{p}_{y}\phi+\hat{z}\hat{p}_{z}\phi\\
&=x\left(\frac{\ii\hbar x}{b^{2}}\exp\left(-\frac{r^{2}}{2b^{2}}\right)\right)\\
&\qquad+y\left(\frac{\ii\hbar y}{b^{2}}\exp\left(-\frac{r^{2}}{2b^{2}}\right)\right)\\
&\qquad+z\left(\frac{\ii\hbar z}{b^{2}}\exp\left(-\frac{r^{2}}{2b^{2}}\right)\right)\\
&=\left(\frac{\ii\hbar}{b^{2}}\exp\left(-\frac{r^{2}}{2b^{2}}\right)\right)(x^{2}+y^{2}+z^{2})\\
&=\frac{\ii\hbar r^{2}}{b^{2}}\exp\left(-\frac{r^{2}}{2b^{2}}\right)
\end{align*}
\[ (\hat{\vect{p}}\cdot\hat{\vect{r}})\phi=\hat{p}_{x}\hat{x}\phi+\hat{p}_{y}\hat{y}\phi+\hat{p}_{z}\hat{z}\phi \]
\begin{align*}
\hat{p}_{x}\hat{x}\phi&=\kakimachigai{\ii\hbar}{-\ii\hbar}\frac{\partial}{\partial x}\left(x\exp\left(-\frac{r^{2}}{2b^{2}}\right)\right)\\
&=-\ii\hbar\biggl(\exp\left(-\frac{r^{2}}{2b^{2}}\right)\\
&\qquad+x\left(-\frac{x}{b^{2}}\exp\left(-\frac{r^{2}}{2b^{2}}\right)\right)\biggr)\\
&=-\ii\hbar\left(1-\frac{x^{2}}{b^{2}}\right)\exp\left(-\frac{r^{2}}{2b^{2}}\right)
\end{align*}
\[ (\hat{\vect{p}}\cdot\hat{\vect{r}})\phi=-\ii\hbar\left(3-\frac{r^{2}}{b^{2}}\right)\exp\left(-\frac{r^{2}}{2b^{2}}\right) \]
\begin{align*}
&(\hat{\vect{p}}\cdot\hat{\vect{r}})\phi-(\hat{\vect{r}}\cdot\hat{\vect{p}})\phi\\
&\quad=-3\ii\hbar\exp\left(-\frac{r^{2}}{2b^{2}}\right)=-3\ii\hbar\phi
\end{align*}
\end{kotaeutsushi}
\begin{maekizuki}{2026-10-14}
\item 交換子の和は $\sum_{j}(\hat{p}_{j}\hat{x}_{j}-\hat{x}_{j}\hat{p}_{j})f$ と成分ごとの差に分けて1行書き，(1) の $\hat{x}\hat{p}_{x}-\hat{p}_{x}\hat{x}=\ii\hbar$ と順序が同じか逆かを見てから符号を決める。
\item 具体的な関数で確かめるときは，(2) と同じ順序の差をとって比べる。順序をそろえていれば (2) の符号の誤りに気付けた。
\end{maekizuki}
\end{mondai}

\level{3}
% toubun: kiroku\toubun\量子力学1\016_2026-10-29.tex
\mondaibangou{7}
\begin{mondai}{016-stationary}{導出}{時間を含むシュレーディンガー方程式と平面波の分散関係，重ね合わせた波の山の速さ（群速度）と定常状態の重ね合わせの時間変化}
\begin{teisuu}
  プランク定数 & $h$ & \qty{6.626e-34}{\joule\second}\\
  換算プランク定数 & $\hbar$ & \qty{1.055e-34}{\joule\second}\\
  電子の質量 & $\melec$ & \qty{9.109e-31}{\kilogram}\\
  電気素量 & $e$ & \qty{1.602e-19}{\coulomb}\\
  真空中の光速 & $c$ & \qty{2.998e8}{\metre\per\second}\\
\end{teisuu}
この問題では，要点整理の時間を含むシュレーディンガー方程式を用いる。1次元で考える。
  \begin{shoumon}
    \item 自由粒子（$V=0$）で $\psi_{k}(x,t)=\exp[\ii(kx-\omega t)]$ が時間を含むシュレーディンガー方程式の解となるための $\omega$ と $k$ の関係を求めよ。また，$E=\hbar\omega$，$p=\hbar k$ を用いてその意味を述べよ。
    \item $0<k_{1}<k_{2}$ とし，$\omega_{1}$，$\omega_{2}$ は (1) の関係を満たすとして，$\psi=\psi_{k_{1}}+\psi_{k_{2}}$ を考える。$\abs{\psi}^{2}$ を $x$ と $t$ の実関数で表し，$\abs{\psi}^{2}$ の山が動く速さ $v_{\mathrm{g}}$ を $k_{1}$，$k_{2}$，$m$，$\hbar$ で表せ。$v_{\mathrm{g}}$ を，2つの平面波に対応する粒子の速さ $\hbar k_{1}/m$，$\hbar k_{2}/m$ と比べよ。さらに，$k_{1}$ と $k_{2}$ が近いとき，$v_{\mathrm{g}}$ は平面波1つの位相速度の何倍かを答えよ。
    \item 幅 $L$ の無限に深い井戸（要点整理）の $\phi_{n}(x)=\sin(n\pi x/L)$ を用いて，$\psi_{n}(x,t)=\phi_{n}(x)\ee^{-\ii E_{n}t/\hbar}$ とおく。$\abs{\psi_{n}}^{2}$ が時間によらないことを示せ。また，$\psi=\psi_{1}+\psi_{2}$ について $\abs{\psi}^{2}$ を求め，$\abs{\psi}^{2}$ が時間変化する角振動数を求めよ。この角振動数を $2\pi$ で割った振動数と，粒子が $n=2$ から $n=1$ の状態に移るときに放出される光子の振動数との関係を述べよ。
  \end{shoumon}
\notekiroku{2026-10-29 B ◎△◎}
\begin{kotaeutsushi}{2026-10-29}
\toi{1}{◎}{}
\[ \ii\hbar\pdv{\psi}{t}=\hat{H}\psi \]
\begin{align*}
\ii\hbar\pdv{\psi_{k}}{t}&=\hbar\omega\exp[\ii(kx-\omega t)]\\
\hat{H}\psi_{k}&=\frac{\hat{p}^{2}}{2m}\psi_{k}\\
&=\frac{1}{2m}\left(-\ii\hbar\pdv{x}\right)^{2}\psi_{k}\\
&=-\frac{\hbar^{2}}{2m}(\ii k)^{2}\exp[\ii(kx-\omega t)]\\
&=\frac{\hbar^{2}k^{2}}{2m}\exp[\ii(kx-\omega t)]
\end{align*}
\[ \hbar\omega=\frac{\hbar^{2}k^{2}}{2m}\qquad E=\frac{p^{2}}{2m} \]
全エネルギーが運動エネルギーのみ。
\toi{2}{△}{方針}
\begin{align*}
\abs{\psi}^{2}&=\left(\psi_{k_{1}}^{*}+\psi_{k_{2}}^{*}\right)\left(\psi_{k_{1}}+\psi_{k_{2}}\right)\\
&=\psi_{k_{1}}^{*}\psi_{k_{1}}+\psi_{k_{2}}^{*}\psi_{k_{2}}\\
&\qquad+\psi_{k_{1}}^{*}\psi_{k_{2}}+\psi_{k_{2}}^{*}\psi_{k_{1}}\\
&=2+\exp[\ii((k_{2}-k_{1})x-(\omega_{2}-\omega_{1})t)]\\
&\qquad+\exp[\ii((k_{1}-k_{2})x-(\omega_{1}-\omega_{2})t)]
\end{align*}
\naze{問いは $x$ と $t$ の実関数で表すことである。2つの指数関数は互いに複素共役なので，$\ee^{\ii\theta}+\ee^{-\ii\theta}=2\cos\theta$ で $\abs{\psi}^{2}=2+2\cos[(k_{2}-k_{1})x-(\omega_{2}-\omega_{1})t]$ までまとめる。}
\begin{align*}
v_{\mathrm{g}}&=\ayamari{\frac{\dd(\omega_{2}-\omega_{1})}{\dd(k_{2}-k_{1})}}{\frac{\omega_{2}-\omega_{1}}{k_{2}-k_{1}}}\\
&=\frac{\dd}{\dd(k_{2}-k_{1})}\left(\frac{\hbar k_{2}^{2}}{2m}-\frac{\hbar k_{1}^{2}}{2m}\right)\\
&=\frac{\dd}{\dd(k_{2}-k_{1})}\,\frac{\hbar}{2m}(k_{2}-k_{1})(k_{2}+k_{1})\\
&=\frac{\hbar}{2m}(k_{1}+k_{2})
\end{align*}
\naze{山は $\cos$ の中身（位相の差）が $2\pi$ の整数倍の点で，その点は $x=[(\omega_{2}-\omega_{1})/(k_{2}-k_{1})]\,t+\text{定数}$ と動く。山の速さは差の比で，微分ではない。以降も微分の形で書いているが，計算は実際には割り算になっていて結果は正しい。}
$\hbar k_{1}/m$ と $\hbar k_{2}/m$ の平均。$k_{1}=k_{2}$ のとき $v_{\mathrm{g}}=\hbar k/m$，\ayamari{1倍}{約2倍}。
\naze{平面波1つの位相速度は (1) の $\omega=\hbar k^{2}/(2m)$ から $\omega/k=\hbar k/(2m)$ で，$\hbar k/m$ ではない（$\hbar k/m$ は粒子の速さ）。$k_{1}\approx k_{2}\approx k$ のとき $v_{\mathrm{g}}\approx\hbar k/m=2\,\omega/k$ で，位相速度の約2倍である。}
\toi{3}{◎}{}
\begin{align*}
\abs{\psi_{n}}^{2}&=\phi_{n}^{*}\ee^{\ii E_{n}t/\hbar}\phi_{n}\ee^{-\ii E_{n}t/\hbar}\\
&=\phi_{n}^{*}\phi_{n}\\
&=\sin^{2}\frac{n\pi x}{L}
\end{align*}
より $t$ に依存しない。
\begin{align*}
\abs{\psi}^{2}&=\abs{\psi_{1}}^{2}+\abs{\psi_{2}}^{2}+\psi_{1}^{*}\psi_{2}+\psi_{2}^{*}\psi_{1}\\
&=\sin^{2}\frac{\pi x}{L}+\sin^{2}\frac{2\pi x}{L}+\psi_{1}^{*}\psi_{2}+\psi_{2}^{*}\psi_{1}
\end{align*}
\begin{align*}
\psi_{1}^{*}\psi_{2}&=\phi_{1}^{*}\phi_{2}\ee^{-\ii(E_{2}-E_{1})t/\hbar}\\
\psi_{2}^{*}\psi_{1}&=\phi_{2}^{*}\phi_{1}\ee^{-\ii(E_{1}-E_{2})t/\hbar}
\end{align*}
\begin{align*}
\abs{\psi}^{2}&=\sin^{2}\frac{\pi x}{L}+\sin^{2}\frac{2\pi x}{L}+\sin\frac{\pi x}{L}\sin\frac{2\pi x}{L}\\
&\qquad\times\left(\ee^{-\ii(E_{2}-E_{1})t/\hbar}+\ee^{-\ii(E_{1}-E_{2})t/\hbar}\right)
\end{align*}
\[ \frac{\abs{E_{2}-E_{1}}}{\hbar} \]
\[ \frac{\abs{E_{2}-E_{1}}}{h}=\frac{\Delta E}{h}=\nu \]
振動数と一致。
\end{kotaeutsushi}
\begin{kizuki}
\item 位相速度は平面波1つの位相 $kx-\omega t$ が一定の点の速さ $\omega/k$ で，自由粒子では $\hbar k/(2m)$，粒子の速さ $\hbar k/m$ の半分である。「何倍か」と問われたら，(1) の $\omega(k)$ から $\omega/k$ を1行で計算してから比べる。
\item \mae{第2章 問題1・10-14} 山の速さは「位相（の差）$=$ 一定」の式を書き，$x$ を $t$ で表して出す。2つの平面波では差の比 $(\omega_{2}-\omega_{1})/(k_{2}-k_{1})$ で，$k_{2}\to k_{1}$ の極限が群速度 $\dv{\omega}{k}$ になる。前回も位相一定の根拠を書き落としたので，速さを出す前にまずこの1行を書く。
\item 「実関数で表せ」と問われたら，複素共役どうしの和を $\ee^{\ii\theta}+\ee^{-\ii\theta}=2\cos\theta$ でまとめる。
\end{kizuki}
\begin{tadashii}
\houshin 平面波を時間を含むシュレーディンガー方程式に代入して $\omega$ と $k$ の関係を得る。重ね合わせの $\abs{\psi}^{2}$ は交差項に位相の差が残るので，その位相の差が一定の点の動きから山の速さを求める。定常状態の重ね合わせでは，交差項にエネルギーの差に比例する角振動数が残る。
\keisan
\shoumonkai{1} $\psi_{k}$ を代入する。
\begin{align*}
  \text{左辺}&=\ii\hbar\pdv{\psi_{k}}{t}=\ii\hbar(-\ii\omega)\psi_{k}=\hbar\omega\,\psi_{k},\\
  \text{右辺}&=-\frac{\hbar^{2}}{2m}\pdv[2]{\psi_{k}}{x}=\frac{\hbar^{2}k^{2}}{2m}\psi_{k}.
\end{align*}
$\psi_{k}\neq0$ なので $\hbar\omega=\hbar^{2}k^{2}/(2m)$，すなわち $\omega=\hbar k^{2}/(2m)$。$E=\hbar\omega$，$p=\hbar k$ を使うと $E=p^{2}/(2m)$ で，自由粒子のエネルギーと運動量の関係を表す。

\shoumonkai{2} $\theta_{j}=k_{j}x-\omega_{j}t$ とおくと $\psi=\ee^{\ii\theta_{1}}+\ee^{\ii\theta_{2}}$ で，
\[ \abs{\psi}^{2}=1+1+\ee^{\ii(\theta_{2}-\theta_{1})}+\ee^{-\ii(\theta_{2}-\theta_{1})}=2+2\cos\left[(k_{2}-k_{1})x-(\omega_{2}-\omega_{1})t\right] \]
山は位相の差 $\theta_{2}-\theta_{1}$ が $2\pi$ の整数倍の点で，その点は $x=(\omega_{2}-\omega_{1})t/(k_{2}-k_{1})+\text{定数}$ と動く。(1) の関係を使うと
\[ v_{\mathrm{g}}=\frac{\omega_{2}-\omega_{1}}{k_{2}-k_{1}}=\frac{\hbar(k_{2}^{2}-k_{1}^{2})}{2m(k_{2}-k_{1})}=\frac{\hbar(k_{1}+k_{2})}{2m} \]
これは2つの粒子の速さ $\hbar k_{1}/m$ と $\hbar k_{2}/m$ の平均である。$k_{1}\approx k_{2}\approx k$ のとき $v_{\mathrm{g}}\approx\hbar k/m$ で，これは要点整理の群速度 $\dv{\omega}{k}=\hbar k/m$ に当たる。平面波1つの位相速度は $\omega/k=\hbar k/(2m)$ なので，$v_{\mathrm{g}}$ はその約2倍である。

\shoumonkai{3} $\theta$ が実数なら $\abs{\ee^{-\ii\theta}}=1$ なので，$\abs{\psi_{n}}^{2}=\phi_{n}^{2}=\sin^{2}(n\pi x/L)$ で時間によらない。$\psi=\psi_{1}+\psi_{2}$ では，$\phi_{n}$ が実数であることを使って
\begin{align*}
  \abs{\psi}^{2}&=\phi_{1}^{2}+\phi_{2}^{2}+\phi_{1}\phi_{2}\left(\ee^{-\ii(E_{2}-E_{1})t/\hbar}+\ee^{\ii(E_{2}-E_{1})t/\hbar}\right)\\
  &=\sin^{2}\frac{\pi x}{L}+\sin^{2}\frac{2\pi x}{L}+2\sin\frac{\pi x}{L}\sin\frac{2\pi x}{L}\cos\frac{(E_{2}-E_{1})t}{\hbar}.
\end{align*}
$\abs{\psi}^{2}$ は角振動数 $\Omega=(E_{2}-E_{1})/\hbar$ で時間変化する。振動数は $\Omega/(2\pi)=(E_{2}-E_{1})/h$ で，$n=2$ から $n=1$ への遷移で放出される光子の振動数 $\nu$（$h\nu=E_{2}-E_{1}$）と一致する。
\kotae{\shoumonkai{1} $\omega=\frac{\hbar k^{2}}{2m}$（$E=p^{2}/(2m)$ を表す） \shoumonkai{2} $\abs{\psi}^{2}=2+2\cos[(k_{2}-k_{1})x-(\omega_{2}-\omega_{1})t]$，$v_{\mathrm{g}}=\frac{\hbar(k_{1}+k_{2})}{2m}$（2つの粒子の速さの平均），位相速度の約2倍 \shoumonkai{3} $\abs{\psi_{n}}^{2}=\sin^{2}(n\pi x/L)$；$\abs{\psi}^{2}=\sin^{2}(\pi x/L)+\sin^{2}(2\pi x/L)+2\sin(\pi x/L)\sin(2\pi x/L)\cos[(E_{2}-E_{1})t/\hbar]$，角振動数 $(E_{2}-E_{1})/\hbar$，振動数は光子の振動数 $(E_{2}-E_{1})/h$ に等しい}
\begin{ronten}
  \item \shoumonkai{1} $\psi_{k}\neq0$ なので両辺の係数を比べてよいことと，得た関係を $E=\hbar\omega$，$p=\hbar k$ で書き直すと自由粒子のエネルギーと運動量の関係になることを述べる。
  \item \shoumonkai{2} $\abs{\psi}^{2}$ の山は位相の差 $\theta_{2}-\theta_{1}$ が $2\pi$ の整数倍の点であり，その点の動く速さを $v_{\mathrm{g}}$ としたことを書く。
  \item \shoumonkai{3} $E_{n}t/\hbar$ が実数なので位相因子の絶対値が1であることと，$\phi_{n}$ が実数であることを使って交差項を $\cos$ にまとめたことを書く。
  \item \shoumonkai{3} 光子の振動数と結びつける根拠として，放出される光子のエネルギーが2つの準位のエネルギーの差に等しいこと（エネルギー保存）を述べる。
\end{ronten}
\kenzan 次元：$\hbar k^{2}/m$ は $(\unit{\kilogram\metre\squared\per\second})(\unit{\per\metre\squared})/\unit{\kilogram}=\unit{\per\second}$。$k_{2}\to k_{1}$ で $v_{\mathrm{g}}\to\hbar k_{1}/m$ となり，粒子の速さに一致する。$\abs{\psi}^{2}$ は0から4の間の値をとる。\shoumonkai{3} で時間平均をとると交差項が消え，$\phi_{1}^{2}+\phi_{2}^{2}$ が残る。$x=L/2$ では $\phi_{2}=0$ なので，交差項はそこで振動しない。
\tsumazuki 定常状態1つの位相因子 $\ee^{-\ii E_{n}t/\hbar}$ は $\abs{\psi_{n}}^{2}$ に現れない。時間変化が見えるのは，エネルギーの異なる状態を重ね合わせたときだけである。
\end{tadashii}
\end{mondai}

\noteowari
\end{document}
